\documentclass[11pt]{article}

\PassOptionsToPackage{table}{xcolor}
\usepackage[final]{acl}
\usepackage{times}
\usepackage{latexsym}
\usepackage[T1]{fontenc}
\usepackage[utf8]{inputenc}
\usepackage{microtype}
\usepackage{inconsolata}
\usepackage{graphicx}
\usepackage{booktabs}
\usepackage{amsmath}
\usepackage{multirow}
\usepackage{enumitem}

\setlength\titlebox{5cm}

\hypersetup{
  pdftitle={A Scaled-Up Empirical Study of Syntactic Language Models},
  pdfauthor={Pengcheng Wang, Zhexuan Zhang, Chao Lou, and Kewei Tu}
}

% Anonymized for double-blind review
\title{A Scaled-Up Empirical Study of Syntactic Language Models}

\author{
  Pengcheng Wang\textsuperscript{1,2} \quad
  Zhexuan Zhang\textsuperscript{1,2} \quad
  Chao Lou\textsuperscript{3} \quad·
  Kewei Tu\textsuperscript{1,2}\thanks{Corresponding author.}
  \\
  \textsuperscript{1}School of Information Science and Technology, ShanghaiTech University
  \\
  \textsuperscript{2}Shanghai Engineering Research Center of Intelligent Vision and Imaging
  \\
  \textsuperscript{3}Independent researcher
  \\
  \texttt{\{wangpch2024, zhangzhx2023, tukw\}@shanghaitech.edu.cn}
  \\
  \texttt{louchaobrn@outlook.com}
}

\begin{document}
\maketitle

\begin{abstract}
Syntactic language models (SLMs) that generate sentences along with their syntactic tree structures have shown promise in both syntactic generalization and downstream tasks. %%Some of these models enforce hierarchical attention patterns using structural attention masks. But two questions remain: whether the benefit comes from tree-structured input, structural attention masks, or both; and whether the approach scales to larger models and diverse tasks.
Prior studies, however, have mostly evaluated smaller models trained on less pretraining data and tested on narrower task suites. In addition, the design space of tree linearization and attention masks remains underexplored.
We conduct a large-scale empirical study of 12 controlled configurations---one terminal baseline and 11 SLM variants spanning different linearizations and masks---at scales up to 1B parameters and 100B training tokens, evaluated across 11 downstream tasks and two syntactic benchmarks.
The experimental results show that SLMs with standard causal attention match or exceed structural-mask variants across downstream and syntactic benchmarks. They require no architectural modifications and remain compatible with FlashAttention.
%The benefit comes from the input format, not the mask.
% We scale our method up with larger parameter sizes and diverse corpora, evaluating on more downstream tasks.
% At the 1B scale, our method achieves the highest 11-task average.
We also observe that when scaling up to the 1B-parameter scale, the top-performing SLM variant no longer shows a clear advantage over a compute-matched non-structural baseline, suggesting that SLMs may become less competitive with scale. Our code is available at \url{https://github.com/qlwpc/TG-Interpolation}.
\end{abstract}

\section{Introduction}

Standard causal language models process text as flat token sequences. Although they can learn distributional proxies for syntax~\cite{hu-etal-2020-systematic}, hierarchical constituency structures---phrases composing into sentences---are left implicit. %%Providing explicit syntactic structure during pretraining could improve syntactic generalization and downstream reasoning. \cite{sartran-etal-2022-transformer}
Syntactic language models (SLMs)~\cite{sartran-etal-2022-transformer, hu-etal-2024-generative, murty-etal-2023-pushdown, zhao-etal-2024-dependency, huang2026GiLT} inject syntactic biases into Transformers~\cite{vaswani2017attention} using a simple technique: jointly modeling surface sentences and linearized versions of their syntactic parse trees.
%%Prior work on SLMs has explored several directions.
%%Transformer Grammars~\cite{sartran-etal-2022-transformer} serialize constituency trees into token sequences and add structural attention masks, allowing closing non-terminals to compose their child constituents. They demonstrate syntactic gains at $\sim$100M scale, but don't explain whether the benefit comes from the tree-structured input tokens, the structural attention mask, or both.
These models typically employ \emph{structural} attention masks that allow a non-terminal to compose its child tokens once they are generated by focusing its attention on them.
\citet{zhao-etal-2025-systematic} conduct a systematic empirical comparison of various design choices for SLMs, but they do not fully explore the complete design space of tree linearization and attention masks.
%%linearization and composition mechanisms.
%%However, their comparison is restricted to composition strategies \textit{within} structural-mask models, and do not compare against baseline models with an equivalent amount of computation.
Moreover, prior SLM studies have mostly evaluated models up to 300M parameters, trained on up to 9B tokens, and tested on only two or three tasks.

In this paper, we conduct a larger-scale empirical study of SLMs, evaluating SLMs with up to 1B parameters, up to 100B training tokens, a contemporary suite of 11 downstream tasks and two syntactic benchmarks. We also investigate new designs of tree linearization and attention masks in a controlled design space of 12 configurations: one terminal baseline and 11 SLM variants.
% In this paper, we address three unanswered questions:
% \begin{enumerate}
%     \item \textbf{Disentanglement}: Does the benefit come from tree-structured input, the structural attention mask, or both? Prior work conflates them.
%     \item \textbf{Scalability and practicality}: Do SLMs remain beneficial at larger scales and on diverse corpora? % Can we drop the inefficient structural masks?
%     \item \textbf{Explicit vs.\ implicit}: How does explicit tree structure compare to implicit structure at matched training budgets?
% \end{enumerate}
% %
% \begin{description}
%     \item[Structural Attention] Tree input with STACK / COMPOSE attention masks (TG) or Nomask. This paradigm requires custom kernel to accelerate attention computing.
%     \item[Causal Attention] Tree input with standard causal attention.
%     \item[Hybrid Attention] Tree input with head-level mask mixing--- splitting half the attention heads using one structure, and the other half using a different structure. This paradigm also requires custom kernel to accelerate attention computing.
%     %\item[Latent-Structure:] We also address the hypothesis that the advantages of SLMs over vanilla LMs come from extra computation resulted from additional input tokens. Specifically, we compare SLMs with vanilla LMs that have dummy tokens inserted into the input sequence ~\cite{goyal2024think}.% Flat input interleaved with pause tokens and causal attention (Pause-1, Pause-2)~\cite{goyal2024think}.%
% \end{description}
%
%%We systematically compare \textbf{12 variants} spanning three attention mechanisms and several tree linearizations variants.
Our findings are as follows: (i) SLMs with standard causal attention match or surpass structural-mask variants across downstream and syntactic benchmarks. They require no architectural modifications and remain compatible with FlashAttention. (ii) At the 1B parameter scale, the top-performing SLM variant no longer shows a clear advantage over a compute-matched non-structural baseline, suggesting that SLMs may become less competitive with scale.

% TODO
In summary, our contributions include:
\begin{itemize}
    \item We scale SLM evaluation to 1B parameters and 100B training tokens, across 11 downstream tasks and two syntactic benchmarks. %and introducing terminal-format evaluation for comparison.
    \item We systematically examine a controlled family of 12 configurations (one terminal baseline and 11 SLM variants) with varying tree linearization and attention mask strategies, disentangling their impact on model performance.
    \item We include compute-matched comparisons between SLMs and non-structural baselines.
    \item Our experiments produce interesting observations regarding SLM design and the competitiveness of SLMs at scale.
\end{itemize}

\section{Related Work}

% \subsection{Syntactic Language Models}

Prior work on syntactic language models (SLMs) has explored a broad range of ways to incorporate linguistic structure into neural sequence modeling. Transformer Grammars (TG)~\cite{sartran-etal-2022-transformer} encode constituency trees as a sequence of actions and enforce hierarchical attention patterns. A subsequent study extends the idea to dependency trees~\cite{zhao-etal-2024-dependency}. Building on this line of work, \citet{zhao-etal-2025-systematic} provide the most systematic comparison of SLM design choices including those of tree linearization and composition mechanisms. These studies show that syntactic inductive bias can improve language modeling and syntactic evaluation, but they typically couple modeling of linearized trees with specialized attention mechanisms.

Other studies incorporate syntactic inductive bias through alternative architectural mechanisms rather than explicit representations. These approaches include latent structure induction during pretraining~\cite{hu-etal-2024-generative}, stack-based recursive state tracking~\cite{murty-etal-2023-pushdown}, and dependency-graph-informed attention mechanisms~\cite{huang2026GiLT}.

% \subsection{Pause Tokens and Latent Computation}

% A separate line of work improves language models by giving them more computation. Pause tokens~\cite{goyal2024think} append learnable tokens to the input, delaying output prediction.
% We include them as a strong baseline with comparable computation with SLMs.
% Scaling Latent Reasoning~\cite{zhu2025scaling} extends pause-token logic to ``looped'' LMs where the same block iterates multiple times per token.
% This positions pause tokens within a broader trend toward latent computation.

% \subsection{Understanding Why Structure Helps}

% \citet{sartran-etal-2022-transformer} argue that structure helps by providing syntactic supervision: non-terminals constrain plausible next words, while TG-style restricted attention promotes composed subtree representations. Such structure improves syntactic generalization, but may trade off against lexical copying and long-context memory.

% % Psycholinguistic probes provide evidence that language models can encode syntactic state. \citet{futrell-etal-2019-neural} examine LMs as psycholinguistic subjects and find evidence for representations of syntactic state. More recent structural-priming evidence from multilingual language models further supports the view on LLMs~\cite{michaelov-etal-2023-structural}.

% This interpretation is also consistent with probability-based syntactic evaluation: \citet{hu-etal-2026-string} show that language-model string probability correlates with grammaticality judgments, supporting the use of log-probability as a syntactic signal. This motivates our use of document-level PPL and terminal-format scoring.

\section{Variants of Syntactic Language Models}

\subsection{Syntactic Language Models}

Let $\mathbf{x} = (x_1, \ldots, x_N)$ be a sentence of $N$ terminal tokens, and let $\mathbf{y}$ be a labeled constituency parse tree for $\mathbf{x}$, where each node carries a phrase category (e.g., \texttt{S}, \texttt{NP}, or \texttt{VP}) and spans a contiguous subsequence of $\mathbf{x}$.
A syntactic language model defines a joint distribution $p(\mathbf{x}, \mathbf{y})$, factorized autoregressively over a linearized token sequence $\mathbf{z} = (z_1, \ldots, z_T)$:
\begin{equation}
p(\mathbf{z} := \textsc{Linearize}(\mathbf{x}, \mathbf{y})) = \prod_{t=1}^{T} p(z_t \mid z_{<t})
\label{eq:joint}
\end{equation}
where $\textsc{Linearize}$ maps the sentence--tree pair into a flat sequence. Under this formulation, prior SLMs can be understood as different ways of instantiating the same autoregressive modeling problem, varying in the \textbf{linearization} that serializes the tree $\mathbf{y}$ into $\mathbf{z}$, and the \textbf{architecture} that parameterizes the conditional probabilities $p(z_t \mid z_{<t})$. In Transformer-based SLMs, this architectural choice is largely realized through the \textbf{attention mechanism}, which governs how syntactic relationships among tokens in
$\mathbf{z}$ are processed.

% Crucially, while Eq.~\ref{eq:joint} defines the modeling objective, existing work differs in \textit{which tokens} the model predicts.
% When $\mathbf{z}$ contains structural tokens, the model may be trained to predict all tokens in $\mathbf{z}$ (including ONT and CNT).


%\paragraph{Latent-Structure} (flat input + learned ``thinking'' tokens). Rather than providing tree structure explicitly, the model receives a flat sequence $\mathbf{z} = (x_1, p_1, \ldots, p_K, x_2, p_1, \ldots, p_K, x_3, \ldots)$ where $\{p_k\}$ are learnable pause tokens~\cite{goyal2024think}.
%The model learns implicit structure through these extra positions with standard causal attention.
%Structure is learned rather than provided, at the cost of additional computation.


% Table~\ref{tab:paradigms} summarizes the four paradigms.

% \begin{table}[t]
%   \centering
%   \resizebox{\columnwidth}{!}{%
%   \begin{tabular}{lccc}
%     \toprule
%     Paradigm & Linearization & Attention & Fast Attn \\
%     \midrule
%     Mask-Structure & Tree-structured & Structural & No \\
%     Sequence-Structure & Tree-structured & Causal & Yes \\
%     Hybrid-Structure & Tree-structured & Hybrid & No \\
%     Latent-Structure & Flat + pause & Causal & Yes \\
%     \bottomrule
%   \end{tabular}%
%   }
%   \caption{Four paradigms as choices of linearization and attention.}
%   \label{tab:paradigms}
% \end{table}

\subsection{Tree Linearization}
\label{sec:tree_linearization}



\begin{table*}[t]
\centering
\renewcommand{\arraystretch}{1.2}
\resizebox{0.94\textwidth}{!}{%
\setlength{\tabcolsep}{4pt}
\begin{tabular}{lcccp{9cm}}
\toprule

\textbf{Name} & \textbf{Left($\mathbf{v}$)} & \textbf{Right($\mathbf{v}$)} & \textbf{Post-process} & \textbf{Example} \\
\midrule
\rowcolor{gray!15}
\multicolumn{5}{l}{\textbf{Linearization}} \\
$\textsc{Lin}^1$ & ONT($v$) & CNT($v$) & - & \texttt{(S Cats (VP chase mice VP) S)} \\
$\textsc{Lin}^2$ & ONT($v$) & 2$\times$CNT($v$) & - & \texttt{(S Cats (VP chase mice VP) VP) S) S)}  \\
$\textsc{Lin}^3$ & ONT($v$) & 3$\times$CNT($v$) & - & \texttt{(S Cats (VP chase mice VP) VP) VP) S) S) S)}  \\
$\textsc{Lin}^1_{\text{--ont}}$ & $\emptyset$ & CNT($v$) & - & \texttt{Cats chase mice VP) S)}  \\
$\textsc{Lin}^1_{\text{--merge}}$ & ONT($v$) & CNT($v$) & Merge CNTs & \texttt{(S Cats (VP chase mice )} \\\midrule
\rowcolor{gray!15}
\multicolumn{5}{l}{\textbf{Non-structural}} \\
$\textsc{Lin}^0$ & $\emptyset$ & $\emptyset$ & - & \texttt{Cats chase mice} \\
$\textsc{Lin}^1_{\text{shuf}}$  &  ONT($v$) & CNT($v$) & Shuffle & \texttt{(S VP) Cats (VP chase S) mice} \\
Pause-1 & $\emptyset$ & $\emptyset$ & Insert 1 <PAUSE> & \texttt{Cats <PAUSE> chase <PAUSE> mice} \\
Pause-2 & $\emptyset$ & $\emptyset$ & Insert 2 <PAUSE> & \texttt{Cats <PAUSE> <PAUSE> chase <PAUSE> <PAUSE> mice} \\
\bottomrule
\end{tabular}%
}
\caption{Overview of linearizations and non-structural token sequences used in this paper. ``Post-process'' refers to operations applied after linearization. ONT and CNT denote opening (e.g., \texttt{(S} and \texttt{(VP}) and closing non-terminal tokens (e.g., \texttt{S)} and \texttt{VP)}), respectively. $v$ indicates that the emitted token carries a node label. ``Merged CNT'' signifies that consecutive CNTs (ignoring labels) are collapsed into a single \texttt{)} token. Finally, $\textsc{Lin}^1_{\text{shuf}}$ is generated by shuffling the $\textsc{Lin}^1$ token sequence while preserving the original terminal order.}
\label{tab:linearizations}
\end{table*}


The function $\textsc{Linearize}(\mathbf{x}, \mathbf{y})$ embeds sentence $\mathbf{x}$ and tree $\mathbf{y}$ into token sequence $\mathbf{z}$. We explore several variants of linearization, all of which share a simple recursive form:
\begin{align}
    \textsc{Lin}(v) &:= \textsc{Left}(v) \notag \\
    &\quad \circ [\, \textsc{Lin}(c) \text{ for } c ~\textsc{in} ~\textsc{children}(v) \,] \notag \\
    &\quad \circ \textsc{Right}(v) \\[1ex]
    \textsc{Lin}(t) &:= t
\end{align}
where $v\in \mathbf{y}$ is a node, $t\in\mathbf{x}$ is a terminal token, and $\circ$ denotes concatenation. $\textsc{Left}(v)$ and $\textsc{Right}(v)$ produce \emph{structural tokens} to enclose the linearized children. $\textsc{children}(v)$ returns the ordered children of $v$ in the tree.

Table~\ref{tab:linearizations} summarizes all linearization variants that we study. $\textsc{Lin}^1$, $\textsc{Lin}^2$ and $\textsc{Lin}^3$ are three variants of linearization, differing only in the number of structural tokens emitted by $\textsc{Right}(v)$. Specifically, $\textsc{Lin}^1$ generates sequences resembling the widely used bracket notation for constituency trees, but with additionally labeled closing brackets. $\textsc{Lin}^2$ was proposed by \citet{sartran-etal-2022-transformer} to define their structural attention mask (referred to as STACK/COMPOSE in our paper). $\textsc{Lin}^3$ is a variant designed to test whether providing more closure tokens helps. We also include two additional variants, $\textsc{Lin}^1_{\text{--ont}}$ and $\textsc{Lin}^1_{\text{--merge}}$, to enable a more systematic comparison. %We select these variants to methodically probe the core design axes of syntax-aware linearization: structural token budget, opening tokens, and per-constituent closure.
Other variations (e.g., different traversal orders) are not included because they would introduce confounds that distract from our central goal: disentangling tree linearization from the attention mask design. %characterizing emergent adaptive computation.

Different linearizations introduce varying numbers of structural tokens, inherently altering the total sequence length for training and inference (Table~\ref{tab:token-budgets}). %For instance, $\textsc{Lin}^1$ and $\textsc{Lin}^2$ produce $2.4\times$ and $3.2\times$ more tokens per raw document, respectively, compared to the baseline $\textsc{Lin}^0$.
To distinguish performance gains due to syntactic structure from gains due simply to increased computation, we include three token-budget-matched baselines alongside the naive terminal sequence. First, $\textsc{Lin}^1_{\text{shuf}}$ shuffles the structural tokens to serve as a control for the necessity of valid structural information. The other two, Pause-1 and Pause-2, insert one or two copies of a dedicated learned \texttt{PAUSE} token~\cite{goyal2024think} after each terminal.


% The core sequence-structured linearization, which we call $textsc{Lin}^1$ (single CNT), produces one ONT and one CNT per constituent:
% \begin{multline}
% \textsc{Lin}^1(v) = \text{ONT}(\ell(v)) \; \circ \; \textsc{Lin}^1(\text{children}(v)) \\
% {}\circ\; \text{CNT}(\ell(v))
% \label{eq:lin1}
% \end{multline}

% where $\circ$ denotes concatenation. For the full tree rooted at $v_{\text{root}}$, $\mathbf{z} = \textsc{Lin}^1(v_{\text{root}})$.
% The standard TG approach~\cite{sartran-etal-2022-transformer} uses a doubled-closure variant:

% \begin{multline}
% \textsc{Lin}^2(v) = \text{ONT}(\ell(v)) \; \circ \; \textsc{Lin}^2(\text{children}(v)) \\
% {}\circ\; \text{CNT}_1(\ell(v)) \; \circ \; \text{CNT}_2(\ell(v))
% \label{eq:lin2}
% \end{multline}

% where $\text{CNT}_1$ and $\text{CNT}_2$ are identical token types (duplicated tokens).
% Under mask-structure attention, $\text{CNT}_1$ attends to constituent-internal tokens (stack) while $\text{CNT}_2$ attends to $\text{CNT}_1$ and prior closures (composition).
% Under causal attention ($\textsc{Lin}^2$ with Sequence-Structure), the model must learn to use this duplicated capacity.


\begin{table}[htbp]
  \centering
  % \small
  \setlength{\tabcolsep}{3.5pt}
  \begin{tabular}{lccc}
    \toprule
    Dataset & $\textsc{Lin}^0$ & $\textsc{Lin}^1$ & $\textsc{Lin}^2$ \\
    \midrule
    BBC News & $\sim$10B & $\sim$24.7B & $\sim$32B \\
    FineWeb-Edu 100BT & $\sim$98B & $\sim$233B & $\sim$301B \\
    \bottomrule
  \end{tabular}
  \caption{Training token counts for each linearization. Structural tokens inflate the token budget 2.4--3.2$\times$ per raw document.}
  \label{tab:token-budgets}
\end{table}



% \subsection{Linearization Variants}
% From these primitives, we define the full space of linearization variants studied in this work (Table~\ref{tab:variants}):

% \begin{description}
%     \item[$\textsc{Lin}^0$ (Terminal):] $\mathbf{z} = \mathbf{x}$, no structural tokens. The standard LM baseline.
%     \item[$\textsc{Lin}^1$ (Tree):] Eq.~\ref{eq:lin1}. One CNT per constituent.
%     \item[$\textsc{Lin}^2$ (TGTree):] Eq.~\ref{eq:lin2}. Two identical CNTs per constituent.
%     \item[$\textsc{Lin}^3$ (Tree-TripleCNT):] Three CNTs per constituent; tests whether more closure tokens help.
%     \item[$\textsc{Lin}^1_{\text{--ont}}$ (Tree-NoONT):] $\textsc{Lin}^1$ without ONT tokens, ablating opening non-terminals.
%     \item[$\textsc{Lin}^1_{\text{--merge}}$ (Tree-Compress):] Consecutive identical CNTs merged into one, ablating per-constituent closure.
%     \item[$\textsc{Lin}^1_{\text{shuf}}$ (Tree-Shuffle):] Tree-linearization tokens randomly permuted.
% \end{description}

\begin{table*}[t]
  \centering
  % \small
  \setlength{\tabcolsep}{6pt}
  \begin{tabular}{cllll}
    \toprule
    \# & Variant & Linearization & Attention  \\
    \midrule
    1 & Terminal & $\textsc{Lin}^0$ & Causal  \\
    2 & Tree & $\textsc{Lin}^1$ & Causal \\
    3 & TGTree & $\textsc{Lin}^2$ & Causal  \\
    4 & TG & $\textsc{Lin}^2$ & STACK/COMPOSE  \\
    5 & TGNomask & $\textsc{Lin}^2$ & Nomask  \\
    6 & TGNomask-Aug$^*$ & $\textsc{Lin}^2$ & Nomask  \\
    7 & Tree-NoONT & $\textsc{Lin}^1_{\text{--ont}}$ & Causal \\
    8 & Tree-Compress & $\textsc{Lin}^1_{\text{--merge}}$ & Causal \\
    9 & Tree-TripleCNT & $\textsc{Lin}^3$ & Causal  \\
    10 & Tree-Shuffle & $\textsc{Lin}^1_{\text{shuf}}$ & Causal   \\
    11 & TGNomask-Mix-TG & $\textsc{Lin}^2$ & Mixing heads (Nomask + STACK/COMPOSE)  \\
    12 & TGTree-Mix-TG & $\textsc{Lin}^2$ & Mixing heads (Causal   + STACK/COMPOSE)  \\
    \bottomrule
  \end{tabular}
  \caption{The 12 configurations in our controlled design space: one terminal baseline and 11 SLM variants. ONT = opening non-terminal, CNT = closing non-terminal. ``Hybrid'' denotes mixing-head models. $^*$TGNomask-Aug differs from TGNomask: CNT$_2$ remains visible to subsequent attention in Aug, while TGNomask prevents subsequent tokens from attending to it.}
  \label{tab:variants}
\end{table*}

\subsection{Structural Attention}
\label{sec:structural_attn}

A Transformer-based SLM can also incorporate structural information directly into the model itself via specially-designed attention mechanisms. We study three attention mechanisms as listed below.

%Structural mask variants (TG, TGNomask, TGNomask-Aug, ...) use $\textsc{Lin}^2$ as their linearization but differ in attention---STACK/COMPOSE (TG), \textbf{Nomask} (TGNomask, TGNomask-Aug), or head-level mask mixing.

\paragraph{Causal Attention} The joint $p(\mathbf{x}, \mathbf{y})$ is modeled through standard autoregressive factorization over $\mathbf{z} = \textsc{Linearize}(\mathbf{x}, \mathbf{y})$ with \textit{no} modification to the attention mechanism.
Constituency structures are encoded entirely in the token sequence; the model learns to route information through Opening Non-Terminal (ONT) and Closing Non-Terminal (CNT) using standard causal attention.
%This paradigm requires no architectural changes relative to standard LMs.

% \paragraph{Structural Attention} The structural attention mechanism factorizes $p(\mathbf{x}, \mathbf{y})$ autoregressively over $z = \textsc{Lin}^2(\mathbf{x}, \mathbf{y})$ and replaces standard causal attention with input-dependent structural masks.
% The original TG approach~\cite{sartran-etal-2022-transformer} uses STACK/COMPOSE attention. % a CNT$_1$ token aggregates all child tokens within the closing constituent via a stack operation; CNT$_2$ (a duplicate of CNT$_1$) implements recursive composition.
% \citet{zhao-etal-2025-systematic} systematically compare structural masks within this paradigm and find that SLMs without sub-constituent masks perform better. We follow their setting as \textbf{Nomask}.
\paragraph{Structural Attention}
Structural attention factorizes $p(\mathbf{x}, \mathbf{y})$ autoregressively over the linearized tree sequence $z = \textsc{Lin}^2(\mathbf{x}, \mathbf{y})$ and uses structure-dependent masks. In the original Transformer Grammar (TG) formulation~\cite{sartran-etal-2022-transformer}, each CNT is duplicated into two actions. CNT$_1$ performs \textsc{COMPOSE} attention over only the current constituent's children, producing a composed representation. CNT$_2$ performs \textsc{STACK} attention over only the representations on the active parse stack. TG then hides the child representations from later positions, forcing them to access the constituent through CNT$_1$. \citet{zhao-etal-2025-systematic} isolate this post-composition masking and show that removing it yields consistently better performance. We call the resulting variant \textbf{Nomask}: it retains the \textsc{COMPOSE} mask at CNT$_1$ but gives all other positions causal access to earlier terminals, ONTs, and CNT$_1$ representations. Thus, ``no mask'' refers specifically to removing TG's post-composition mask, not to unrestricted bidirectional attention. Our default Nomask still hides CNT$_2$ states from later positions; Nomask-Aug exposes them as well.

\paragraph{Hybrid Attention} The hybrid attention mechanism factorizes $p(\mathbf{x}, \mathbf{y})$ autoregressively over $z = \textsc{Lin}^2(\mathbf{x}, \mathbf{y})$ and replaces standard causal attention with \textbf{head-level mask mixing}: rather than applying a single mask type across all attention heads, different heads receive different masks. Hybrid attention enables the model to combine multiple structural attention patterns within a single Transformer layer.
Since enumerating all possible head-level mask assignments is computationally infeasible, we restrict ourselves to a straightforward setting: half of the attention heads receive one type of structural mask, and the remaining heads use another (either causal or structural). This configuration provides a tractable way to study hybrid attention.

Table~\ref{tab:variants} summarizes the 12 controlled configurations across the tree-linearization and attention-mask designs studied in our framework.

\section{Terminal Format Evaluation}
\label{sec:terminal-format-evaluation}

To evaluate SLMs on multiple-choice (MC) tasks, we follow the OLMES completion/cloze formulation~\cite{gu-etal-2025-olmes} with multi-shot in-context demonstrations.
Each answer is treated as a candidate completion, rather than as a label predicted from a displayed option list.
We parse both the context and continuation with Benepar~\cite{kitaev-klein-2018-constituency} and use their respective 1-best constituency parses.
Appendix~\ref{sec:olmes-eval} provides the complete input construction.

SLM continuations interleave terminal content with structural ONT and CNT tokens, creating a choice between scoring the complete sequence and scoring only terminal tokens.
Let $\mathbf{p}$ denote the evaluation context and $\mathbf{c}_j$ the linearized continuation for candidate $j$.
Let $\mathcal{T}(\mathbf{c}_j)$ and $\mathcal{N}(\mathbf{c}_j)$ denote its terminal and non-terminal indices, respectively.
The full score includes both token types, whereas the terminal score includes only answer-content tokens:
\begin{align}
S_{\text{full}}(\mathbf{c}_j \mid \mathbf{p}) &= \sum_{t \in \mathcal{T} \cup \mathcal{N}} \log p(c_{j,t} \mid \mathbf{p}, \mathbf{c}_{j,<t}) \\
S_{\text{term}}(\mathbf{c}_j \mid \mathbf{p}) &= \sum_{t \in \mathcal{T}} \log p(c_{j,t} \mid \mathbf{p}, \mathbf{c}_{j,<t})
\label{eq:decom}
\end{align}

We compare these scoring rules on a held-out eight-task validation suite (Validation-8; 17,691 examples).
Terminal scoring increased macro-average accuracy by 2.16 points for Tree (95\% CI [1.35, 2.99]) and by 0.62 points for TGTree (95\% CI [$-$0.09, 1.35]).
Based on this result, we use terminal scoring as the primary metric and predict $\arg\max_j S_{\mathrm{term}}(\mathbf{c}_j \mid \mathbf{p})$ for MC tasks.
Appendix~\ref{sec:decom} defines the validation suite and reports the corresponding full-score sensitivity analysis.


\section{Experiments}

\subsection{Setup}

We train all models from scratch for one epoch using the same OLMo-based~\cite{groeneveld-etal-2024-olmo} decoder-only Transformer family across scales. Learning rates and batch sizes follow the Step Law~\cite{li2025predictable}. Causal attention models use FlashAttention-2~\cite{dao2023flashattention2}, and structural-mask models require FlexAttention~\cite{dong2024flex}. Further details of model architectures and hyperparameter selection are provided in Appendix~\ref{sec:training}.

Table~\ref{tab:variants} defines the controlled design space used for our linearization and attention-mask ablations. To address the varying training token budgets induced by different linearizations, we compare SLMs with pause-token baselines~\cite{goyal2024think}: $\textsc{Pause-}K = (x_1, \underbrace{p, \ldots, p}_{K}, x_2, \underbrace{p, \ldots, p}_{K}, x_3, \ldots)$, where $p$ is a dedicated learned \texttt{PAUSE} token. For $K=1$, the computational load is roughly comparable to $\textsc{Lin}^1$ ($2.4\times$); for $K=2$, it is roughly comparable to $\textsc{Lin}^2$ ($3.2\times$).

We additionally include a 100M-parameter \emph{architectural reference baseline} outside this design space. Pushdown~\cite{murty-etal-2023-pushdown} jointly trains the language model and an attachment parser, whose predictions induce the stack-depth attention biases used at inference. We reimplement Pushdown in our OLMo-based Transformer architecture and pretrain it for one epoch on silver parse trees from the same BBC News corpus. This reproduction ensures that Pushdown uses the same Transformer backbone and pretraining data as the other models. We include it for reference only: it is not counted among the 12 controlled configurations or used in the ablations and scaling selection.

\subsection{Evaluation Methods}

\paragraph{Document-Level Language Modeling.}
\label{sec:doc-ppl}
Our SLMs define a joint distribution $p(\mathbf{x}, \mathbf{y})$ over a sentence and its constituency tree (Eq.~\ref{eq:joint}).
For each sentence--tree pair, we form $\mathbf{z}=\textsc{Linearize}(\mathbf{x},\mathbf{y})$ and compute its joint probability autoregressively as $p(\mathbf{x},\mathbf{y})=\prod_{t=1}^{|\mathbf{z}|}p(z_t\mid\mathbf{z}_{<t})$.
The sentence probability $p(\mathbf{x})=\sum_{\mathbf{y}\in\mathcal{Y}_{\mathbf{x}}}p(\mathbf{x},\mathbf{y})$ requires marginalizing over all valid parses and is therefore intractable.
Following \citet{sartran-etal-2022-transformer} and \citet{zhao-etal-2025-systematic}, we approximate each conditional marginal using 300 labeled proposal trees from a CRF constituency parser~\cite{kitaev-klein-2018-constituency}:
\begin{align}
\widetilde p(\mathbf{x}_i \mid \mathbf{h}_i)
&= \sum_{\mathbf{y} \in \mathcal{Y}'_i}
p(\mathbf{x}_i,\mathbf{y} \mid \mathbf{h}_i).
\label{eq:marginal}
\end{align}
\noindent Here, $\mathcal{Y}'_i\subset\mathcal{Y}_{\mathbf{x}_i}$ is the proposal set, and $K=|\mathcal{Y}'_i|=300$ is its size. The truncated sum is a lower bound on the full marginal and therefore yields an upper bound on perplexity.

For document-level evaluation, marginalization over the parse histories of all preceding sentences would incur exponential cost.
We therefore follow the greedy single-path approximation of prior work.
For each preceding sentence $j$, we retain the proposal with the highest joint probability under the evaluated model,
\begin{equation}
\hat{\mathbf{y}}_j=\arg\max_{\mathbf{y}\in\mathcal{Y}'_j}p(\mathbf{x}_j,\mathbf{y}\mid\mathbf{h}_j).
\end{equation}
The selected parses form the prefix $\mathbf{h}_i=\textsc{Linearize}(\mathbf{x}_{<i},\hat{\mathbf{y}}_{<i})$, and only the current sentence is marginalized over all 300 proposals.

For a corpus of $D$ documents, let $\mathbf{X}_d=(\mathbf{x}_{d,1},\ldots,\mathbf{x}_{d,N_d})$ denote document $d$, where $\mathbf{x}_{d,i}$ is its $i$-th sentence.
Document-level perplexity is then
\begin{align}
\mathcal{L}_K
&= \sum_{d=1}^{D}\sum_{i=1}^{N_d}
\log \widetilde p(\mathbf{x}_{d,i}\mid\mathbf{h}_{d,i}), \nonumber\\
\text{PPL}_{\text{syn}} &= \exp(-\mathcal{L}_K/M).
\label{eq:doc-ppl}
\end{align}
\noindent Here, $M$ is the total number of scored terminal tokens. Except for the model-specific procedures in Appendix~\ref{sec:ppl-protocol-details}, tree models use Eq.~\ref{eq:marginal}; flat and pause-token models use terminal-token autoregressive PPL.
%for $\textsc{Lin}^0$ (terminal), there is no tree structure to marginalize over ($\mathbf{y}$ is vacuous) and Eq.~\ref{eq:doc-ppl} reduces to standard autoregressive perplexity.
%For pause-token models ($\textsc{Lin}^{\text{pause}}_K$), there is similarly no parse tree---the model defines $p(\mathbf{x})$ directly over the flat sequence $\mathbf{z} = (\mathbf{x}, p_1, \ldots, p_K)$, and Eq.~\ref{eq:doc-ppl} reduces to standard perplexity computed over terminal tokens.

% \paragraph{Syntactic evaluations.} We probe syntactic competence with two benchmarks. SG~\cite{hu-etal-2020-systematic} tests six targeted generalization subtasks, for which we decode with the word-synchronous beam search method of \citet{stern-etal-2017-effective}; BLiMP~\cite{warstadt-etal-2020-blimp-benchmark} evaluates whether models assign higher likelihood to grammatical sentences in minimal pairs across 12 grammatical phenomena.

% \paragraph{Downstream evaluations.}
\paragraph{Evaluations.} We measure downstream performance with task-specific fine-tuning on XSum summarization~\cite{narayan-etal-2018-dont}, reporting ROUGE-1/2/L~\cite{lin-2004-rouge}, and on BoolQ~\cite{clark-etal-2019-boolq} for boolean question answering. For each fixed BBC pretrained checkpoint, we run five independent fine-tuning seeds and report mean $\pm$ sample standard deviation for both tasks. For FineWeb-Edu models, we adopt the OLMES task formats and completion/cloze formulation~\cite{gu-etal-2025-olmes}. Our 11-task suite contains WinoGrande~\cite{sakaguchi2020winogrande}, HellaSwag~\cite{zellers-etal-2019-hellaswag}, BoolQ, MMLU~\cite{hendrycks2021measuring}, MMLU-Redux~\cite{gema-etal-2025-done}, OpenBookQA~\cite{mihaylov-etal-2018-suit}, CommonsenseQA~\cite{talmor-etal-2019-commonsenseqa}, SocialIQA~\cite{sap-etal-2019-social}, ARC-Easy and ARC-Challenge~\cite{clark2018think}, and PIQA~\cite{bisk2020piqa}. We additionally evaluate syntactic competence on SG~\cite{hu-etal-2020-systematic} and BLiMP~\cite{warstadt-etal-2020-blimp-benchmark}.

\begin{table*}[htbp]
  \centering
  \footnotesize
  \setlength{\tabcolsep}{2.2pt}
  \begin{tabular}{lccc@{\hspace{6pt}}c@{\hspace{6pt}}c@{\hspace{6pt}}c@{\hspace{6pt}}c}
    \toprule
    Model & Linearization & Attention & XSum$\uparrow$ & BoolQ$\uparrow$ & SG$\uparrow$ & BLiMP$\uparrow$ & Doc-PPL$^\ddagger$$\downarrow$ \\
    \midrule
    Terminal-100M & $\textsc{Lin}^0$ & Causal & $22.02 \pm 0.06$ & $66.70 \pm 0.36$ & 69.80 & 70.77 & 9.89 \\
    Tree-100M & $\textsc{Lin}^1$ & Causal & $22.84 \pm 0.03$ & $68.07 \pm 0.35$ & 81.40 & 80.95 & 10.01 \\
    TGTree-100M & $\textsc{Lin}^2$ & Causal & \textbf{23.03 $\pm$ 0.02} & $68.12 \pm 0.62$ & 79.47 & 81.46 & 10.21 \\
    TG-100M & $\textsc{Lin}^2$ & S/C & $20.15 \pm 0.04$ & $64.50 \pm 0.40$ & 79.19 & 83.28 & 11.02 \\
    TGNomask-100M & $\textsc{Lin}^2$ & Nomask & $21.91 \pm 0.04$ & $66.66 \pm 0.45$ & 78.62 & 81.78 & 9.97 \\
    TGNomask-Aug-100M & $\textsc{Lin}^2$ & Nomask & $22.12 \pm 0.02$ & $67.01 \pm 0.28$ & 78.68 & \textbf{83.63} & 10.08 \\
    Tree-NoONT-100M & $\textsc{Lin}^1_{\text{--ont}}$ & Causal & $22.41 \pm 0.04$ & \textbf{68.67 $\pm$ 0.32} & \textbf{85.38} & 81.00 & 9.95 \\
    Tree-Compress-100M & $\textsc{Lin}^1_{\text{--merge}}$ & Causal & $22.70 \pm 0.05$ & $68.48 \pm 0.49$ & 78.27 & 81.40 & 9.91 \\
    Tree-TripleCNT-100M & $\textsc{Lin}^3$ & Causal & $22.43 \pm 0.04$ & $68.09 \pm 0.12$ & 81.57 & 83.42 & 10.26 \\
    Tree-Shuffle-100M & $\textsc{Lin}^1_{\text{shuf}}$ & Causal & $21.56 \pm 0.09$ & $68.08 \pm 0.33$ & 76.56 & 67.77 & 13.41 \\
    TGNomask-Mix-TG-100M & $\textsc{Lin}^2$ & N+S/C & $21.87 \pm 0.04$ & $66.13 \pm 0.54$ & 75.77 & 83.14 & 10.01 \\
    TGTree-Mix-TG-100M & $\textsc{Lin}^2$ & C+S/C & $22.16 \pm 0.07$ & $67.91 \pm 0.36$ & 78.84 & 82.75 & 10.00 \\
    Pause-1-100M & Pause-1 & Causal & $22.38 \pm 0.05$ & $62.04 \pm 0.19$ & 75.11 & 70.85 & \textbf{9.77} \\
    Pause-2-100M & Pause-2 & Causal & $22.25 \pm 0.06$ & $65.70 \pm 3.45$ & 76.97 & 72.59 & 10.32 \\
    \midrule
    \rowcolor{gray!15}
    Pushdown-100M & $\textsc{Lin}^0$ & Stack bias & $21.06 \pm 0.05$ & $65.54 \pm 0.29$ & 74.86 & 74.97 & 13.29 \\
    \midrule
    Terminal-500M & $\textsc{Lin}^0$ & Causal & $20.75 \pm 0.06$ & $69.69 \pm 0.60$ & \textbf{71.11} & 64.74 & \textbf{2.83} \\
    Tree-500M & $\textsc{Lin}^1$ & Causal & \textbf{22.25 $\pm$ 0.05} & $70.72 \pm 0.35$ & 63.13 & 76.09 & 3.18\\
    TGTree-500M & $\textsc{Lin}^2$ & Causal & $22.01 \pm 0.08$ & $70.63 \pm 0.32$ & 65.95 & \textbf{77.65} & 3.28\\
    TGNomask-Aug-500M & $\textsc{Lin}^2$ & Nomask & $21.77 \pm 0.06$ & \textbf{70.78 $\pm$ 0.74} & 57.58 & 75.97 & 3.19\\
    \bottomrule
  \end{tabular}
  \caption{Key 100M and 500M results on BBC News. The shaded row is an architectural reference baseline for reference only, not a member of the controlled variant family. XSum reports R-AVG; XSum and BoolQ show mean $\pm$ sample standard deviation over five fine-tuning seeds. Other metrics are single evaluations. S/C = STACK/COMPOSE; N+S/C and C+S/C mix it at head level with Nomask and causal attention, respectively. Stack bias is Pushdown's stack-depth bias, derived from its attachment parser. Highest means per scale and best available single-evaluation values are in \textbf{bold}. $\ddagger$Doc-PPL protocols are detailed in Appendix~\ref{sec:ppl-protocol-details}.}
  \label{tab:exp1-results}
\end{table*}

\subsection{Experiment 1: Disentangling Tree Linearization and Attention}
\label{sec:exp1}

For training corpus construction, we use the Benepar parser~\cite{kitaev-klein-2018-constituency}, which is trained on news-domain data, to obtain constituency parses for BBC News articles filtered from each subset of the FineWeb dataset~\cite{penedo2024fineweb}.
The resulting corpus covers approximately 15.2M articles from BBC News domains and contains $\sim$10B terminal tokens.
We then train all 100M controlled configurations and comparison baselines on this corpus for one epoch and evaluate them on XSum ROUGE, BoolQ, SG, and BLiMP.
Our goal is to identify which tree linearization and which attention mechanisms drive the syntactic benefit, and to screen members of the controlled design space for scaling.
Table~\ref{tab:exp1-results} separates the shaded architectural reference baseline from the controlled configurations and compute controls.

\paragraph{Structural Attention vs.\ Causal Attention.}
SLMs with causal attention outperform structural attention models on downstream tasks.
On XSum, TGTree and Tree have the highest means (23.03 and 22.84), exceeding TGNomask-Aug (22.12), Terminal (22.02), and TG (20.15); BoolQ shows the same ordering of means among these models.
The full STACK/COMPOSE mask (TG) underperforms even the terminal baseline on both tasks, while a softer variant (TGNomask-Aug) yields a partial recovery.
Structural attention models exhibit a clear syntax--utility tradeoff: TG and its variants lead only on BLiMP (up to 83.63), where the hard constraints of STACK/COMPOSE attention enforce grammatical consistency, but the same rigidity penalizes general language modeling.
Causal attention models suffer no such tradeoff---TGTree has the highest 100M XSum mean, while Tree combines a 22.84 XSum mean with 81.40 SG; TGNomask-Aug manages only 78.68 SG.
The mask is neither necessary nor beneficial for most objectives. %; the tree input format alone suffices.

\paragraph{Architectural Reference Baseline.}
Pushdown does not improve XSum or BoolQ over Terminal but improves both syntactic scores (74.86 and 74.97) under its model-specific inference procedure.

\paragraph{Hybrid Attention Fails to Improve.}
Hybrid attention variants (TGNomask-Mix-TG, TGTree-Mix-TG) yield no benefit over causal Tree SLM, adding complexity without improving performance.

\paragraph{Pause-token Controls.}
The pause-token baselines improve over Terminal on XSum (22.38 and 22.25 vs.\ 22.02), indicating that additional latent-computation positions can help generation.
Nevertheless, the compute-matched tree models remain stronger across downstream and syntactic evaluations: Tree exceeds Pause-1 on XSum, BoolQ, SG, and BLiMP, and TGTree likewise exceeds Pause-2 on all four metrics.
Thus, the gains from tree linearization cannot be explained solely by the larger token budget.

\paragraph{Document-Level Perplexity.}
%The tree-marginalized document-level PPL (Eq.~\ref{eq:doc-ppl}) provides an orthogonal signal that tests pure language modeling quality independent of task-specific decoding.
%At 100M (full results in Appendix~\ref{sec:full-100m}),
The PPL of Tree remains close to that of Terminal (10.01 vs.\ 9.89), with a gap of 0.12 PPL points, indicating that linearizing trees into the input sequence incurs only a small LM cost.
The STACK/COMPOSE mask is the exception: TG degrades PPL to 11.02, consistent with its downstream underperformance.
TGNomask is competitive with causal tree models, supporting \citeposs{zhao-etal-2025-systematic} finding that removing sub-constituent masks improves LM quality.
Among the tree-linearization variants, Tree-Compress achieves the lowest PPL at 9.91. Tree-Shuffle obtains 13.41 under its terminal-only protocol (Appendix~\ref{sec:ppl-protocol-details}), so we exclude it from direct comparisons with CRF-marginalized PPL values.
Among all 100M models, Pause-1 achieves the lowest PPL at 9.77, slightly outperforming Terminal at 9.89, whereas Pause-2 reaches 10.32.
Except for Tree-Shuffle, the reproduced PPL values cluster near Terminal even when downstream scores differ, indicating that PPL and downstream performance provide complementary signals.

% \paragraph{Tree Linearization Ablations} Within all SLMs with causal attention, Opening non-terminals primarily aid generation (higher XSum with ONTs), not syntactic discrimination (SG is higher without them), suggesting they help the model anticipate upcoming constituent structure.
% Tree-Shuffle causes a substantial BLiMP decline.
% Merging consecutive identical CNTs (Tree-Compress) reduces SG, suggesting that fine-grained closure information is useful.

\paragraph{Tree Linearization Ablations. }
The experiments on all SLMs with causal attention form an ablation study of tree linearization.
Tree-NoONT improves syntactic discrimination but impairs generation, suggesting ONTs primarily help the model anticipate upcoming constituent structures rather than enforce syntactic constraints. Tree-Shuffle performs substantially worse on BLiMP. Doubling CNTs (TGTree) achieves the best trade-off between syntax and generation; tripling CNTs (Tree-TripleCNT) raises SG relative to TGTree (81.57 vs.\ 79.47) at the cost of generation quality and reaches 83.42 on BLiMP. Merging consecutive identical closing non-terminals, as in Tree-Compress, reduces SG, suggesting that fine-grained closure information helps targeted syntactic generalization.

\paragraph{Scaling Up.} After the 100M experiments, we select four variants and train them at 500M parameters on BBC News to verify that the advantage of tree linearization with causal attention persists before committing to larger-scale experiments.
Table~\ref{tab:exp1-results} shows the results, confirming that Tree and TGTree maintain their downstream lead over the terminal baseline and the structural attention variant.
SG scores decline for all three scaled tree-input models, while Terminal changes little; we discuss this pattern in \S\ref{sec:sg-discussion}.
% We proceed to 1B with Tree and TGTree as the primary syntactic variants.

There is a practical note on throughput: SLMs with causal attention benefit from FlashAttention compatibility, whereas SLMs with structural attention require slower custom attention implementations.
Table~\ref{tab:efficiency} reports measured training throughput.

\begin{table}[tbp]
  \centering
  \fontsize{10pt}{12pt}\selectfont
  \setlength{\tabcolsep}{6pt}
  \begin{tabular}{lccc}
    \toprule
    Scale & Causal Mask & Structural Mask & Ratio \\
    \midrule
    100M & 62,175 token/s & 37,975 token/s & 1.6$\times$ \\
    500M & 12,550 token/s & 8,316 token/s & 1.5$\times$ \\
    \bottomrule
  \end{tabular}
  \caption{Training throughput on H800 GPUs. Causal models use FlashAttention-2, whereas structural-mask models use FlexAttention. The mixing-head mask achieves 3,375 token/s at the 100M scale, which is 18$\times$ slower than causal attention.}
  \label{tab:efficiency}
\end{table}

The throughput gap persists at scale, and mixing-head mask models are 18$\times$ slower.
Combined with the performance evidence above, the causal tree approach is the clear practical choice for scaling.

Based on these results, we select tree linearization Tree ($\textsc{Lin}^1$) and TGTree ($\textsc{Lin}^2$), both with causal attention, as the primary variants for scaling, dropping all mixing-head and structural attention approaches.

\subsection{Experiment 2: Scaling to 1B on FineWeb-Edu --- Main Result}
\label{sec:exp2}

FineWeb-Edu-100BT~\cite{penedo2024fineweb} is a $\sim$100B-token sampled subset of the FineWeb-Edu corpus, filtered for educational quality and spanning diverse domains.
We train five models at 1B parameters on this data: Terminal, Tree, TGTree, Pause-1 and Pause-2.
Evaluation uses the 11-task OLMES suite~\cite{gu-etal-2025-olmes}, with the content-focused terminal score as the primary MC metric (\S\ref{sec:terminal-format-evaluation}), together with the SG and BLiMP syntactic benchmarks. Table~\ref{tab:main-result} reports per-task downstream accuracy, and Table~\ref{tab:1b-syntax} reports the syntactic evaluation results. We quantify evaluation-set uncertainty using 10,000 paired bootstrap resamples; full confidence intervals and tests appear in Appendix~\ref{sec:bootstrap}, and Appendix~\ref{sec:decom} reports sensitivity to full versus terminal scoring.

\begin{table*}[htbp]
  \centering
  %\small
  \setlength{\tabcolsep}{1.5pt}
  \begin{tabular}{lccccccccccccc}
    \toprule
    Model & WG$\uparrow$ & HS$\uparrow$ & BQ$\uparrow$ & MM$\uparrow$ & MMR$\uparrow$ & OBQA$\uparrow$ & CSQA$\uparrow$ & SIQA$\uparrow$ & ARCE$\uparrow$ & ARCC$\uparrow$ & PIQA$\uparrow$ & \textbf{AVG}$\uparrow$ \\
    \midrule
    Terminal & 58.17 & 61.53 & 58.65 & 36.79 & 37.61 & 43.60 & 53.24 & 48.98 & 69.30 & 43.48 & 73.88 & 53.20 \\
    Tree & 61.64 & 54.10 & 63.85 & 36.88 & 37.53 & 44.00 & 56.02 & 48.46 & 68.42 & 45.15 & 73.07 & 53.56 \\
    TGTree & \textbf{63.38} & 63.51 & \textbf{67.55} & 37.76 & 38.42 & \textbf{45.80} & \textbf{58.15} & 50.31 & 70.00 & 44.82 & 74.97 & \textbf{55.88} \\
    Pause-1 & 58.96 & 63.05 & 65.41 & 37.70 & 38.81 & 42.80 & 56.18 & \textbf{51.28} & \textbf{74.04} & 43.48 & \textbf{75.14} & 55.17 \\
    Pause-2 & 62.59 & \textbf{64.45} & 63.39 & \textbf{38.38} & \textbf{39.46} & 44.00 & 56.76 & 48.67 & 72.98 & \textbf{48.16} & 74.97 & 55.80 \\
    \bottomrule
  \end{tabular}
  \caption{Primary content-focused terminal-score accuracy at 1B on FineWeb-Edu, based on the scoring definition in Section~\ref{sec:terminal-format-evaluation}. WG = WinoGrande, HS = HellaSwag, BQ = BoolQ, MM = MMLU, MMR = MMLU-Redux, OBQA = OpenBookQA, CSQA = CommonsenseQA, SIQA = SocialIQA, ARCE = ARC-Easy, ARCC = ARC-Challenge. Best per task in \textbf{bold}; Table~\ref{tab:bootstrap-ci} reports 95\% bootstrap intervals.}
  \label{tab:main-result}
\end{table*}

\paragraph{Main result.}
TGTree has the highest 11-task point estimate (55.88, +2.68 over Terminal), but the bootstrap ranking is not decisive: its probability of ranking first is 0.29, compared with 0.28 for Pause-2. Excluding BoolQ reverses their point-estimate order (Pause-2: 55.04; TGTree: 54.71). Relative to Terminal, TGTree shows significant improvements on five of the 11 tasks; the full paired comparisons are reported in Appendix~\ref{sec:bootstrap}.

\paragraph{Sensitivity to scoring format.}
As defined in Section~\ref{sec:terminal-format-evaluation}, terminal scoring is our primary content-focused metric; full scoring is a complementary sensitivity analysis. On Validation-8, terminal scoring changes 17.58\% of the Tree rankings and 12.35\% of the TGTree rankings relative to full scoring. It raises macro-average accuracy by 2.16 points for Tree (95\% CI [1.35, 2.99]) and 0.62 points for TGTree (95\% CI [$-$0.09, 1.35]). This supports the primary-metric choice: Tree improves significantly, while TGTree's point estimate rises without a statistically distinguishable change. Appendix~\ref{sec:decom} reports the task-level results and limitations.

\paragraph{Explicit Structure vs.\  Dummy Tokens.}
TGTree and the compute-matched Pause-2 baseline have nearly identical 11-task point estimates (55.88 vs.\ 55.80). In paired tests, TGTree is significantly better only on BoolQ (+4.16 [2.02, 6.30], $p<.001$), while Pause-2 is significantly better on HellaSwag (+0.94 [0.19, 1.68], $p=.014$); the other nine task differences are not significant at $\alpha=.05$. A separate ten-task validation evaluation yields the same aggregate conclusion (57.35 vs.\ 57.09; difference +0.26, 95\% CI [$-$0.68, 1.20]; Table~\ref{tab:validation10}). These models therefore show compute-matched competitive performance with complementary task-level strengths, rather than dominance. A combination of explicit structure and learned latent computation is a natural direction for future work.

\begin{table}[tbp]
  \centering
  \setlength{\tabcolsep}{8pt}
  \begin{tabular}{lcc}
    \toprule
    Model & BLiMP$\uparrow$ & SG$\uparrow$ \\
    \midrule
    Terminal & 79.65 & 88.01 \\
    Tree & \textbf{81.79} & \textbf{89.95} \\
    TGTree & 80.98 & 89.00 \\
    Pause-1 & 80.83 & 86.54 \\
    Pause-2 & 80.36 & 86.59 \\
    \bottomrule
  \end{tabular}
  \caption{Syntactic evaluation at 1B on FineWeb-Edu. For Tree and TGTree, SG uses word-synchronous beam search with a maximum linearized decoding length of six times the terminal sequence length.}
  \label{tab:1b-syntax}
\end{table}

\paragraph{Syntactic evaluation at 1B.}
Tree has the highest point estimate on both syntactic benchmarks, exceeding Terminal by 2.14 points on BLiMP and 1.94 points on SG. TGTree also exceeds Terminal on both benchmarks, whereas the compute-matched pause-token baselines improve BLiMP but not SG. Without pretraining replicates, we treat these results as task-level evidence rather than a general model ordering.

\subsection{Scaling Trend Summary}

\begin{table*}[htbp]
  \centering
  \setlength{\tabcolsep}{12pt}
  \begin{tabular}{lllcc}
    \toprule
    Scale & Data & Metric & Terminal & Best SLM ($\Delta$) \\
    \midrule
    100M & BBC & XSum R-AVG & $22.02 \pm 0.06$ & TGTree: $23.03 \pm 0.02$ (+1.01) \\
    500M & BBC & XSum R-AVG & $20.75 \pm 0.06$ & Tree: $22.25 \pm 0.05$ (+1.50) \\
    1B & BBC & XSum R-AVG & $20.47 \pm 0.06$ & Tree: $21.42 \pm 0.08$ (+0.94) \\
    1B & FW-Edu & 11-task AVG & 53.20 & TGTree: 55.88 (+2.68) \\
    \bottomrule
  \end{tabular}
  \caption{Scaling trend across model sizes and datasets. BBC XSum values are mean $\pm$ sample standard deviation over five fine-tuning seeds; the FineWeb-Edu row is a single evaluation. $\Delta$ is the difference between means for BBC and the point-estimate difference for FineWeb-Edu.}
  \label{tab:scaling-summary}
\end{table*}

Table~\ref{tab:scaling-summary} shows positive point-estimate differences over Terminal across all scales and datasets.
On BBC, the mean gap between Terminal and the best evaluated SLM is non-monotonic (+1.01 $\rightarrow$ +1.50 $\rightarrow$ +0.94; the 1B BBC run was supplementary), possibly reflecting saturation on the narrow 10B-token corpus.
On diverse FineWeb-Edu data, the point-estimate gap is +2.68 on the 11-task average. As discussed in \S\ref{sec:exp2}, however, the identity of the top model is not stable under bootstrap resampling; the stronger evidence is the set of significant task-level improvements over Terminal.

% \begin{table}[htbp]
%   \centering
%   \resizebox{\columnwidth}{!}{%
%   \begin{tabular}{lccc}
%     \toprule
%     Variant & XSum$\uparrow$ & SG$\uparrow$ & BLiMP$\uparrow$ \\
%     \midrule
%     Tree (baseline) & \textbf{22.82} & 81.40 & 80.95  \\
%     Tree-NoONT & 22.29 & \textbf{85.38} & 81.00  \\
%     Tree-Compress & 22.63 & 78.27 & 81.40  \\
%     Tree-Shuffle & 21.76 & 61.59 & 71.68  \\
%     TGTree ($2\times$CNT) & 22.69 & 79.47 & 81.46 \\
%     Tree-TripleCNT ($3\times$CNT) & 21.18 & 81.57 & ---\\
%     \bottomrule
%   \end{tabular}%
%   }
%   \caption{Ablation results at 100M BBC News.}
%   \label{tab:ablations}
% \end{table}

\section{Discussion}

\subsection{Why Does Tree Linearization Help Without Structural Attention?}
\label{sec:why-works}

We propose three complementary hypotheses:

\paragraph{Implicit multi-task learning hypothesis.} The model jointly predicts terminal tokens and syntactic structure, making constituency boundaries explicit through non-terminals instead of inferring them only from co-occurrence statistics. \citet{sartran-etal-2022-transformer} attribute the perplexity gains of such joint modeling to the ``additional syntactic supervision'' that non-terminals provide by constraining admissible word classes at each position.

\paragraph{Compositional bottleneck hypothesis.} CNT positions may act as ``summary points'' that encourage compression of constituent information. Under this hypothesis, causal attention allows the model to learn which information to route through CNT positions, unlike TG's hard-coded COMPOSE attention. This flexibility may be more useful for diverse data and may contribute to the larger gains on FineWeb-Edu. We do not directly test this mechanism.

\paragraph{Adaptive computation hypothesis.} Pause-$K$ controls for extra tokens by allocating a fixed number of learned latent tokens after each terminal. Tree models instead use parse-dependent structural tokens: richer constituency structures introduce more ONT/CNT positions at phrase boundaries. This placement may allocate additional computation near constituent composition. TGTree and Pause-2 are statistically indistinguishable on nine of the 11 downstream tasks, so our results motivate this hypothesis but do not establish it.

%\paragraph{Soft bias $>$ hard constraint.} Causal attention lets the model learn \textit{when} to use syntactic structure; STACK/COMPOSE forces it even when harmful. The finding that TG (hard mask) underperforms on general tasks while excelling at BLiMP supports this: the hard mask over-constrains attention, trading general capability for grammatical sensitivity. Causal tree models get the syntactic benefit without the rigidity penalty.

\subsection{Syntactic Degradation on BBC News}
\label{sec:sg-discussion}

On BBC News, scaling Tree from 100M to 500M lowers SG from 81.40 to 63.13 and BLiMP from 80.95 to 76.09, whereas Terminal SG remains stable (69.80 to 71.11). This decline is not accompanied by worse in-domain modeling: Tree's document PPL improves from 10.01 to 3.18. For parse ranking, we score 300 Benepar n-best parses for each of 100 BBC sentences. At $K=300$, the 500M model selects Benepar's top parse more often than the 100M model (0.81 vs.\ 0.72). Thus, the larger model better fits both BBC News text and the analyses preferred by the news-domain parser, despite weaker generalization to SG and BLiMP.

The degradation does not recur on the more diverse FineWeb-Edu corpus. At 1B, Tree and TGTree score 89.95 and 89.00 on SG, both above Terminal's 88.01. On BLiMP, they score 81.79 and 80.98, again above Terminal's 79.65 (Table~\ref{tab:1b-syntax}). Together, the improved in-domain fit, stronger parse ranking, and cross-corpus contrast are consistent with overfitting to the narrow BBC News distribution. Because corpus diversity, tokenizer, and training scale are not independently controlled, this evidence supports but does not establish a causal domain-overfitting explanation.

\section{Conclusion}

We conduct a scaled-up empirical study of syntactic language models, evaluating one terminal baseline and 11 SLM variants at scales up to 1B parameters and 100B training tokens. Our evaluation spans 11 downstream tasks and two syntactic benchmarks. Our controlled comparisons disentangle tree linearization from attention design: SLMs with standard causal attention match or surpass structural-mask variants across downstream and syntactic benchmarks. They also require no architectural modifications and remain compatible with FlashAttention, making causal tree models the more practical SLM design. At the 1B scale, however, the top-performing SLM and a compute-matched non-structural baseline show competitive aggregate performance with complementary task-level strengths, rather than clear SLM dominance. This finding suggests that SLMs may become less competitive with scale and that their advantages should be assessed against equally compute-intensive non-structural alternatives.


\section*{Limitations}

Constructing linearized tree inputs requires a parser for both training data and inputs. This adds preprocessing cost, and parsing errors may propagate into the training signal. We do not quantify the impact of parse quality. In addition, we investigate only English, leaving syntactic structures in other languages unexplored.

% Our terminal MC score is a content-focused scoring choice, not an exact sentence marginal. It retains one silver parse as context while omitting its structural-token probabilities. The validation results shows that switching formats can change rankings in either direction, so SLM MC results should be interpreted with respect to the reported scoring format.

The compositional bottleneck and adaptive computation accounts are hypotheses motivated by the observed comparisons, not established mechanisms. The core conclusions are supported by controlled comparisons across linearizations, attention masks, pause-token baselines, and scales, but we do not directly intervene on internal representations or attention dynamics to prove how tree tokens are used. Such mechanistic probes are left to future work.

% At every model scale, each pretraining configuration is represented by a single run; the five BBC seeds capture fine-tuning variance only. The paired bootstrap intervals quantify uncertainty due to the finite evaluation sets, but they do not capture variation across pretraining seeds. Moreover, the per-task tests are exploratory and are reported without correction for multiple comparisons. Multi-seed pretraining and confirmatory evaluation on held-out task suites are left to future work.

% Acknowledgments (anonymized)
\section*{Acknowledgments}
This work was supported by the HPC Platform of ShanghaiTech University, the SIST Computing Platform and the Core Facility Platform of Computer Science and Communication, SIST, ShanghaiTech University.


\bibliography{anthology-1,anthology-2,custom}

\appendix

\section{Training Details}
\label{sec:training}

% TODO:
\subsection{Implementation and Hyperparameters}

We use an OLMo-based~\cite{groeneveld-etal-2024-olmo} decoder-only Transformer with SwiGLU~\cite{shazeer2020gluSWIGLU} activation, RoPE~\cite{su2024roformerROPE} positional encoding, RMSNorm~\cite{zhang2019rootRMS} for pre-normalization, and tied input and output embeddings.
Under this architecture, we train different SLMs with 100M, 500M, and 1B parameters. Table~\ref{tab:scales} lists the three scale configurations.

\begin{table}[htbp]
  \centering
  \small
  \setlength{\tabcolsep}{3pt}
  \begin{tabular}{lccccc}
    \toprule
    Scale & $d_\text{model}$ & MLP ratio & Layers & Heads & Params \\
    \midrule
    100M & 768 & 4 & 12 & 12 & $\sim$100M \\
    500M & 1408 & 8 & 16 & 16 & $\sim$500M \\
    1B & 2048 & 8 & 16 & 16 & $\sim$1B \\
    \bottomrule
  \end{tabular}
  \caption{Model configurations at each scale.}
  \label{tab:scales}
\end{table}

We train all models from scratch for one epoch. Training uses the AdamW optimizer~\cite{loshchilov2017decoupledADAM} with $\beta_1{=}0.9$, $\beta_2{=}0.95$, $\epsilon{=}10^{-8}$, weight decay of 0.1, and gradient clipping at norm of 1.0. We use a cosine learning rate scheduler with 2,000 warmup steps; after the learning rate decays to $10^{-5}$, it is held constant for the rest of training. The maximum sequence length is 2,048 tokens. The learning rate $\eta$ and batch size $B$ are set according to the Step Law~\cite{li2025predictable}:

\begin{equation}
    \begin{gathered}
       \eta(N,D) = 1.79N^{-0.713}D^{0.307} \\
B(D) = 0.58D^{0.571}
    \end{gathered}
\label{eq:steplaw}
\end{equation}
where $N$ is the number of non-embedding model parameters and $D$ is the size of training dataset.

\subsection{Computational Costs}
At each parameter scale, we report the aggregate training time across all model variants evaluated at that scale.
The 100M-parameter variants required 133 hours on 4 H800 GPUs.
The 500M-parameter variants required 125 hours on 4 H800 GPUs.
The 1B-parameter variants required 385 hours (approximately 16.0 days) on 64 A800 GPUs.

\subsection{Tokenizer and Structural Tokens}
We use tokenizers with reserved non-terminal tokens (e.g., \texttt{NP)}, \texttt{(S}, and \texttt{VP)}).
The tokenizer is based on the GPT-2 tokenizer~\cite{radford2019language} or the Qwen3 tokenizer~\cite{yang2025qwen3}.
For tree linearization models, opening non-terminals (ONT) and closing non-terminals (CNT) are represented as reserved tokens so that linearized constituency structure is preserved during pretraining and evaluation.

\section{Evaluation Setup Details}

\subsection{Downstream Fine-Tuning on BBC News: XSum and BoolQ}
On BBC News, we evaluate XSum and BoolQ through task-specific fine-tuning from the pretrained checkpoint, with the optimizer state reset.
All models are fine-tuned with AdamW ($\beta_1{=}0.9$, $\beta_2{=}0.95$, weight decay 0.1) and a cosine learning rate schedule with 100 warmup steps decaying to $10^{-6}$.

\paragraph{XSum.} We fine-tune all models for 3 epochs with a peak learning rate of $6\times10^{-5}$ and a global batch size of 40.
During evaluation, we generate summaries with beam search using a beam size of 6 and a maximum length of 150 tokens.
For tree linearization models, beam search proceeds at the word level (word-sync) with a maximum of 75 word steps and a grammar constraint of at most 4 consecutive opening non-terminals.
ROUGE-1, ROUGE-2, and ROUGE-L are computed with the standard evaluate library; R-AVG denotes their arithmetic mean.

\paragraph{BoolQ.} We fine-tune all models for 5 epochs with peak LR $3\times10^{-4}$, global batch size 40.
We measure accuracy on the validation split.

The XSum and BoolQ entries in Table~\ref{tab:exp1-results} are means $\pm$ sample standard deviations over five fine-tuning seeds (42, 2026, 6198, 13171, and 31723). Pushdown fine-tuning uses right-binarized, unary-collapsed gold spans. Its XSum decoder conditions on the gold parse of the source and jointly searches summary tokens and attachment actions (beam size 6, at most 4 consecutive reductions); BoolQ uses tree-aware candidate scoring with beam size 20.

\subsection{Syntactic Generalization (SG)}
\label{sec:sg-eval-setup}
Following \citet{hu-etal-2020-systematic}, we evaluate the six SG test suites without fine-tuning by comparing target-region surprisal~\cite{hale-2001-probabilistic}.
For a tagged span $\mathbf{w}=\mathbf{x}_{s:e}$ and its left context $\mathbf{C}=\mathbf{x}_{<s}$, we define
\[
\begin{aligned}
S(\mathbf{w}\mid\mathbf{C})
  &= -\log p(\mathbf{w}\mid\mathbf{C}) \\
  &= -\log p(\mathbf{x}_{\leq e})+\log p(\mathbf{x}_{<s}).
\end{aligned}
\]
The logarithm base does not affect the SG predictions because the task-specific formulas only compare surprisals.

For terminal and pause-token models, we compute this span surprisal by summing token-level cross-entropy losses over the tagged span in one forward pass.
For tree-linearization models other than Tree-Shuffle, we approximate prefix probabilities with DFS word-synchronous beam search~\cite{stern-etal-2017-effective}.
At each synchronized terminal boundary $t$, the beam $\mathcal{B}_t$ contains the $b$ highest-scoring latent structural prefixes $\mathbf{y}^{(h)}_t$.
We form the beam-truncated marginal and the corresponding span score as
\[
\begin{aligned}
\widetilde{p}(\mathbf{x}_{\leq t})
  &= \sum_{h\in\mathcal{B}_t}
     p(\mathbf{x}_{\leq t},\mathbf{y}^{(h)}_t), \\
\widetilde{S}(\mathbf{w}\mid\mathbf{C})
  &= -\log\widetilde{p}(\mathbf{x}_{\leq e}) \\
  &\quad +\log\widetilde{p}(\mathbf{x}_{<s}).
\end{aligned}
\]
These span surprisals are the scores supplied to the task-specific formulas of \citet{hu-etal-2020-systematic}.

At each terminal step, the decoder must emit the observed next terminal, but may first expand non-terminal hypotheses through depth-first search.
Hypotheses synchronize after emitting the terminal, and the $b$ highest-scoring hypotheses are retained.
Following \citet{murty-etal-2023-pushdown} and \citet{sartran-etal-2022-transformer}, we set $b{=}300$; the within-step sub-beam is 300 and the terminal fast-track size is 30.
All such SG evaluations cap the linearized decoding length at $\max(6L,10)$, where $L$ is the terminal-token sequence length.
No separate limit is imposed on the total or consecutive number of non-terminal tokens.

Tree-Shuffle uses terminal-format teacher-forced scoring at test time. Pushdown instead incrementally marginalizes over attachment decisions with a beam size of 300.

\subsection{BLiMP}
BLiMP~\cite{warstadt-etal-2020-blimp-benchmark} evaluates whether models assign higher likelihood to grammatical sentences in minimal pairs across 12 grammatical phenomena.
For each BLiMP minimal pair evaluated with CRF-based tree marginalization, we obtain the 300 highest-scoring parse proposals ($K=300$) from the Benepar parser~\cite{kitaev-klein-2018-constituency} and compute the tree-marginalized probability of each sentence via $p(\mathbf{x}) \approx \sum_{k=1}^{K} p(\mathbf{x}, \mathbf{y}_k)$, following the same marginalization procedure as document-level PPL (Eq.~\ref{eq:marginal}).
A pair is scored correct if the good sentence receives higher marginalized probability than the bad sentence.
Each of the 67 subtasks contains 1,000 pairs.
For terminal and pause-token models, $K=1$ (no tree marginalization).
Tree-Shuffle likewise uses terminal-format scoring with $K=1$. For Pushdown, each candidate tree supplies a fixed sequence of attachment actions, and the score marginalizes over the 300 candidate trees.

\subsection{Model-Specific Document-PPL Protocols}
\label{sec:ppl-protocol-details}

CRF-parsed tree-linearization models other than Tree-Shuffle use the 300-proposal truncated marginal in Eq.~\ref{eq:marginal}. Under the same greedy parser-selected history, truncation lower-bounds the exact current-sentence probability and therefore upper-bounds its PPL. Tree-Shuffle uses reproduced terminal-only scoring. Pushdown sums joint token--attachment probabilities over up to 300 native attachment candidates for the current sentence and uses candidate~0 for each preceding sentence. Because the Tree-Shuffle and Pushdown procedures differ from CRF-based marginalization, their PPL values are contextual comparisons.

\subsection{OLMES Multiple-Choice Evaluation on FineWeb-Edu}
\label{sec:olmes-eval}
All 11 downstream tasks use the completion/cloze formulation standardized in OLMES~\cite{gu-etal-2025-olmes}. This formulation evaluates each answer string separately as a natural-language completion instead of predicting an answer label. All evaluations are multi-shot. We use 5-shot prompts for HellaSwag, WinoGrande, MMLU, MMLU-Redux, and OpenBookQA, and 3-shot prompts for BoolQ, ARC-Easy, ARC-Challenge, PIQA, SocialIQA, and CommonsenseQA.

We preprocess the natural-language content of every demonstration, target prompt, and answer option with Benepar~\cite{kitaev-klein-2018-constituency}.
For each parsed component $x$, we retain only its 1-best constituency parse $\hat{\mathbf{y}}(x)$ as a silver tree.
Depending on the task template, fields are parsed separately or the completed candidate sequence is parsed and split at the cloze boundary.
Task-specific template text is represented with the same constituency-tree notation.
The resulting trees are converted to the linearization required by each model.
Thus, the syntactic inputs differ in representation but share the same underlying terminal prompt and answer text.

For each demonstration, the parsed gold answer suffix is appended to its parsed prompt.
These demonstrations precede the target instance in every candidate prefix $\mathbf{p}_j$.
We then append the candidate-specific parsed suffix $\mathbf{c}_j$ and score only that suffix conditional on $\mathbf{p}_j$.
The prediction is $\arg\max_j S(\mathbf{c}_j \mid \mathbf{p}_j)$, using the task-specific option-score normalization. The primary results use $S_{\text{term}}$ in Eq.~\ref{eq:decom}; Appendix~\ref{sec:decom} also reports $S_{\text{full}}$.
For WinoGrande, the answer fills the blank in a candidate-specific prefix, while the common post-blank text is the scored suffix, following OLMES.
Unlike BLiMP, this MC protocol uses one silver tree per textual component and uses neither multiple parse proposals nor tree marginalization.

\section{Bootstrap Confidence Intervals and Paired Tests}
\label{sec:bootstrap}

We estimate uncertainty for the 1B FineWeb-Edu downstream evaluation using $B=10{,}000$ bootstrap resamples with seed 2027. For each task, examples are resampled with replacement and the same sampled indices are used for every model. Table~\ref{tab:bootstrap-ci} reports percentile 95\% confidence intervals for accuracy. For the macro-average rows, tasks rather than examples are resampled; these wider intervals characterize sensitivity to task-suite composition. The probability that each model ranks first under this task resampling is 0.09 (Terminal), 0.11 (Tree), 0.29 (TGTree), 0.23 (Pause-1), and 0.28 (Pause-2). Excluding BoolQ changes these probabilities to 0.11, 0.10, 0.26, 0.22, and 0.29, respectively.

\begin{table*}[p]
  \centering
  \footnotesize
  \renewcommand{\arraystretch}{1.05}
  \setlength{\tabcolsep}{2.8pt}
  \begin{tabular}{lccccc}
    \toprule
    Task & Terminal & Tree & TGTree & Pause-1 & Pause-2 \\
    \midrule
    WG   & 58.17 [55.49, 60.93] & 61.64 [58.96, 64.33] & 63.38 [60.69, 66.06] & 58.96 [56.20, 61.64] & 62.59 [59.91, 65.27] \\
    HS   & 61.53 [60.57, 62.48] & 54.10 [53.14, 55.09] & 63.51 [62.57, 64.45] & 63.05 [62.09, 63.99] & 64.45 [63.51, 65.39] \\
    BQ   & 58.65 [56.97, 60.34] & 63.85 [62.23, 65.50] & 67.55 [65.93, 69.11] & 65.41 [63.82, 67.06] & 63.39 [61.71, 65.08] \\
    MM   & 36.79 [35.99, 37.59] & 36.88 [36.06, 37.69] & 37.76 [36.96, 38.56] & 37.70 [36.88, 38.51] & 38.38 [37.57, 39.18] \\
    MMR  & 37.61 [36.37, 38.89] & 37.53 [36.28, 38.77] & 38.42 [37.18, 39.70] & 38.81 [37.56, 40.07] & 39.46 [38.23, 40.72] \\
    OBQA & 43.60 [39.20, 48.00] & 44.00 [39.60, 48.40] & 45.80 [41.40, 50.20] & 42.80 [38.40, 47.00] & 44.00 [39.60, 48.40] \\
    CSQA & 53.24 [50.45, 56.02] & 56.02 [53.32, 58.80] & 58.15 [55.36, 60.93] & 56.18 [53.40, 58.97] & 56.76 [53.97, 59.54] \\
    SIQA & 48.98 [46.83, 51.23] & 48.46 [46.26, 50.67] & 50.31 [48.16, 52.56] & 51.28 [49.03, 53.58] & 48.67 [46.47, 50.92] \\
    ARCE & 69.30 [65.44, 72.98] & 68.42 [64.56, 72.11] & 70.00 [66.14, 73.68] & 74.04 [70.35, 77.54] & 72.98 [69.30, 76.49] \\
    ARCC & 43.48 [38.13, 49.16] & 45.15 [39.46, 50.84] & 44.82 [39.13, 50.50] & 43.48 [38.13, 49.16] & 48.16 [42.47, 53.85] \\
    PIQA & 73.88 [71.87, 75.90] & 73.07 [71.06, 75.08] & 74.97 [72.96, 76.93] & 75.14 [73.18, 77.09] & 74.97 [73.01, 76.93] \\
    \midrule
    AVG$^\dagger$ & 53.20 [46.38, 60.30] & 53.56 [46.79, 60.42] & 55.88 [48.39, 63.08] & 55.17 [47.86, 62.66] & 55.80 [48.68, 63.15] \\
    AVG w/o BQ$^\dagger$ & 52.66 [45.29, 60.26] & 52.53 [45.38, 60.00] & 54.71 [47.18, 62.60] & 54.14 [46.27, 62.17] & 55.04 [47.46, 62.91] \\
    \bottomrule
  \end{tabular}
  \caption{Accuracy point estimates and 95\% bootstrap confidence intervals for the terminal-format 1B evaluation. Per-task rows use paired resampling of evaluation examples. $^\dagger$Macro-average intervals resample tasks and therefore measure task-suite sensitivity rather than uncertainty within a fixed suite. Task abbreviations follow Table~\ref{tab:main-result}.}
  \label{tab:bootstrap-ci}
\end{table*}

Table~\ref{tab:bootstrap-tests} reports the paired contrasts most relevant to our claims: TGTree versus the Terminal baseline and each tree model versus its compute-matched pause-token baseline. We compute two-sided bootstrap $p$-values from the paired difference distribution. The tests are exploratory and uncorrected for multiple comparisons; $p<.05$ is used only as a descriptive threshold.

\begin{table*}[p]
  \centering
  \footnotesize
  \renewcommand{\arraystretch}{1.05}
  \setlength{\tabcolsep}{6pt}
  \begin{tabular}{lccc}
    \toprule
    Task & TGTree $-$ Terminal & Tree $-$ Pause-1 & TGTree $-$ Pause-2 \\
    \midrule
    WG   & +5.21 [2.29, 8.13]; $p<.001$ & +2.68 [$-$0.32, 5.68]; $p=.082$ & +0.78 [$-$2.05, 3.63]; $p=.604$ \\
    HS   & +1.98 [1.22, 2.74]; $p<.001$ & $-$8.95 [$-$9.96, $-$7.89]; $p<.001$ & $-$0.94 [$-$1.68, $-$0.19]; $p=.014$ \\
    BQ   & +8.90 [6.48, 11.31]; $p<.001$ & $-$1.57 [$-$3.27, 0.12]; $p=.074$ & +4.16 [2.02, 6.30]; $p<.001$ \\
    MM   & +0.97 [0.26, 1.67]; $p=.009$ & $-$0.82 [$-$1.53, $-$0.09]; $p=.027$ & $-$0.62 [$-$1.30, 0.08]; $p=.078$ \\
    MMR  & +0.81 [$-$0.32, 1.96]; $p=.172$ & $-$1.27 [$-$2.40, $-$0.18]; $p=.024$ & $-$1.04 [$-$2.14, 0.05]; $p=.069$ \\
    OBQA & +2.16 [$-$1.40, 5.80]; $p=.262$ & +1.21 [$-$2.80, 5.20]; $p=.592$ & +1.76 [$-$1.80, 5.40]; $p=.367$ \\
    CSQA & +4.92 [2.46, 7.45]; $p<.001$ & $-$0.15 [$-$2.46, 2.21]; $p=.934$ & +1.39 [$-$0.90, 3.69]; $p=.248$ \\
    SIQA & +1.34 [$-$0.67, 3.38]; $p=.210$ & $-$2.81 [$-$4.86, $-$0.77]; $p=.007$ & +1.65 [$-$0.36, 3.68]; $p=.120$ \\
    ARCE & +0.67 [$-$3.16, 4.56]; $p=.761$ & $-$5.63 [$-$9.30, $-$1.93]; $p=.002$ & $-$3.00 [$-$6.67, 0.70]; $p=.113$ \\
    ARCC & +1.36 [$-$4.01, 6.69]; $p=.666$ & +1.64 [$-$3.68, 6.69]; $p=.570$ & $-$3.32 [$-$9.03, 2.34]; $p=.261$ \\
    PIQA & +1.07 [$-$0.65, 2.83]; $p=.247$ & $-$2.09 [$-$3.92, $-$0.27]; $p=.024$ & $-$0.02 [$-$1.74, 1.69]; $p=1.000$ \\
    \bottomrule
  \end{tabular}
  \caption{Paired bootstrap differences in percentage points with 95\% confidence intervals and uncorrected two-sided $p$-values. Positive values favor the model named first. Tree--Pause-1 and TGTree--Pause-2 are compute-matched comparisons.}
  \label{tab:bootstrap-tests}
\end{table*}

\section{SG Subtask Breakdown}

Table~\ref{tab:sg-subtasks} breaks down SG performance by subtask for key 100M variants. Agreement shows the largest improvement (Tree +41.4 vs.\ Terminal). Center Embedding shows ceiling effects across all models (87.5--94.6). For Long-Distance Dependencies, Tree and Tree-NoONT improve by 3.21 points and TGTree by 5.96 points, whereas TG decreases by 0.46 points relative to Terminal.

\begin{table*}[p]
  \centering
  \small
  \renewcommand{\arraystretch}{1.08}
  \setlength{\tabcolsep}{12pt}
  \begin{tabular}{lcccccc}
    \toprule
    Model & Agr$\uparrow$ & CtrEmb$\uparrow$ & GrdnPath$\uparrow$ & GrsSyn$\uparrow$ & Lic$\uparrow$ & LDD$\uparrow$ \\
    \midrule
    Terminal & 39.66 & 87.50 & 78.29 & 84.78 & 61.28 & 68.35 \\
    Tree & 81.03 & 89.29 & 84.21 & 83.70 & 80.08 & 71.56 \\
    TGTree & 67.24 & 94.64 & 82.24 & 88.04 & 69.92 & 74.31 \\
    TG & 79.31 & 92.86 & 81.58 & 79.35 & 71.80 & 67.89 \\
    Tree-NoONT & 87.93 & 89.29 & 85.53 & 97.83 & 80.08 & 71.56 \\
    \bottomrule
  \end{tabular}
  \caption{SG subtask breakdown at 100M. Agr = Agreement, CtrEmb = Center Embedding, GrdnPath = Garden Path Effects, GrsSyn = Gross Syntactic Expectation, Lic = Licensing, LDD = Long Distance Dependencies. $\uparrow$ = higher is better.}
  \label{tab:sg-subtasks}
\end{table*}

Table~\ref{tab:sg-subtasks-1b} reports the corresponding breakdown for the 1B FineWeb-Edu models. Tree and TGTree use the maximum linearized decoding length of six times the terminal sequence length described in Appendix~\ref{sec:sg-eval-setup}.

\begin{table*}[t]
  \centering
  \fontsize{8pt}{10pt}\selectfont
  \setlength{\tabcolsep}{4pt}
  \begin{tabular}{lccccccc}
    \toprule
    Model & Agr$\uparrow$ & CtrEmb$\uparrow$ & GrdnPath$\uparrow$ & GrsSyn$\uparrow$ & Lic$\uparrow$ & LDD$\uparrow$ & SG$\uparrow$ \\
    \midrule
    Terminal & \textbf{87.72} & 94.74 & 93.42 & 98.91 & 74.81 & 78.44 & 88.01 \\
    Tree & 85.96 & \textbf{100.00} & \textbf{98.03} & \textbf{100.00} & 78.20 & 77.52 & \textbf{89.95} \\
    TGTree & 80.70 & \textbf{100.00} & 94.08 & \textbf{100.00} & \textbf{78.95} & \textbf{80.28} & 89.00 \\
    Pause-1 & 75.44 & \textbf{100.00} & 94.08 & 98.91 & 73.31 & 77.52 & 86.54 \\
    Pause-2 & 82.46 & 94.74 & 95.39 & \textbf{100.00} & 72.18 & 74.77 & 86.59 \\
    \bottomrule
  \end{tabular}
  \caption{SG subtask breakdown at 1B on FineWeb-Edu. Best per metric in \textbf{bold}; abbreviations follow Table~\ref{tab:sg-subtasks}.}
  \label{tab:sg-subtasks-1b}
\end{table*}


\section{Full-versus-Terminal Score Sensitivity}
\label{sec:decom}

We evaluate the scoring choice on the held-out Validation-8 data by decomposing the MC score into terminal and non-terminal components (Eq.~\ref{eq:decom}). Validation-8 comprises the validation splits of WinoGrande, HellaSwag, OpenBookQA, CommonsenseQA, SocialIQA, ARC-Easy, ARC-Challenge, and PIQA, totaling 17,691 examples. We draw fixed in-context demonstrations from the training data and verify that their normalized text does not overlap with the scored examples. Table~\ref{tab:decom} reports both scores for the 1B Tree and TGTree models.

\begin{table*}[htbp]
  \centering
  \small
  \setlength{\tabcolsep}{5pt}
  \begin{tabular}{lrrr@{\hspace{12pt}}rrr}
    \toprule
    & \multicolumn{3}{c}{Tree} & \multicolumn{3}{c}{TGTree} \\
    \cmidrule(lr){2-4}\cmidrule(lr){5-7}
    Task & Full & Term & $\Delta$ & Full & Term & $\Delta$ \\
    \midrule
    WinoGrande & 61.17 & 61.64 & +0.47 & 64.64 & 63.38 & $-$1.26 \\
    HellaSwag & 37.83 & 43.92 & +6.08 & 47.11 & 49.45 & +2.34 \\
    OpenBookQA & 33.60 & 33.60 & +0.00 & 34.20 & 34.00 & $-$0.20 \\
    CommonsenseQA & 54.79 & 55.86 & +1.06 & 57.17 & 57.49 & +0.33 \\
    SocialIQA & 43.65 & 45.60 & +1.94 & 44.42 & 45.80 & +1.38 \\
    ARC-Easy & 71.05 & 74.74 & +3.68 & 73.68 & 74.39 & +0.70 \\
    ARC-Challenge & 39.80 & 41.81 & +2.01 & 38.80 & 39.80 & +1.00 \\
    PIQA & 70.24 & 72.25 & +2.01 & 72.63 & 73.29 & +0.65 \\
    \midrule
    Macro average & 51.52 & 53.68 & +2.16 & 54.08 & 54.70 & +0.62 \\
    \bottomrule
  \end{tabular}
  \caption{Full and terminal accuracy (\%) on Validation-8; $\Delta=\mathrm{Term}-\mathrm{Full}$. Macro averages are unweighted. Paired 95\% bootstrap intervals for the macro $\Delta$ are [1.35, 2.99] for Tree and [$-$0.09, 1.35] for TGTree (10,000 resamples; seed 2026).}
  \label{tab:decom}
\end{table*}

Non-terminal tokens have higher mean log-probability than terminal tokens by 2.68 nats for Tree and 2.86 nats for TGTree, but this scale difference does not by itself determine option rankings. With terminal scoring, 7.46\% of Tree examples flip to correct and 5.30\% to wrong. The corresponding TGTree rates are 4.92\% and 4.30\%. Other ranking changes switch between two incorrect options and therefore do not affect accuracy. At the macro level, terminal scoring improves Tree by 2.16 points (95\% CI [1.35, 2.99]). TGTree rises by 0.62 points, but its interval includes zero (95\% CI [$-$0.09, 1.35]). Two TGTree tasks favor full scoring. These results support terminal scoring as the primary content-focused metric while motivating full scoring as a complementary sensitivity analysis.

The two scores address related but distinct questions. $S_{\text{full}}$ scores the candidate together with one parser-selected silver tree. $S_{\text{term}}$ accumulates only terminal-token factors while retaining that same tree as context. Neither score is the exact marginal sentence likelihood $\sum_{\mathbf{y}}p(\mathbf{x},\mathbf{y})$. Accordingly, $S_{\text{term}}$ is primary, while $S_{\text{full}}$ measures sensitivity to including the parser-selected structural continuation.

\subsection{Separate Validation-10 Comparison}
\label{sec:validation10}

Validation-10 adds BoolQ and a held-out MMLU validation set to the eight tasks above, yielding 22,207 scored examples. For MMLU, five fixed demonstrations are selected for each of 57 subjects and removed from scoring, leaving 1,246 examples; MMLU accuracy is macro-averaged across subjects. All other task accuracies are micro-averaged, and the headline score is the unweighted macro-average over ten tasks. We use 10,000 paired bootstrap resamples with seed 2026.

\begin{table}[htbp]
  \centering
  \small
  \setlength{\tabcolsep}{6pt}
  \begin{tabular}{lcc}
    \toprule
    Model & Macro accuracy & $\Delta$ vs.\ Terminal (95\% CI) \\
    \midrule
    Terminal & 54.43 & reference \\
    Tree & 54.75 & +0.32 [$-$0.67, 1.30] \\
    TGTree & \textbf{57.35} & +2.92 [1.95, 3.84] \\
    Pause-1 & 56.47 & +2.05 [1.15, 2.93] \\
    Pause-2 & 57.09 & +2.66 [1.82, 3.49] \\
    \bottomrule
  \end{tabular}
  \caption{Terminal-score accuracy (\%) on Validation-10. The paired TGTree--Pause-2 difference is +0.26 points (95\% CI [$-$0.68, 1.20]).}
  \label{tab:validation10}
\end{table}

\end{document}
